\section{The construction side}\label{sec:construction}

Our aim at the outset was a linear-time transposer. We did not find one. This section records what the
results above require of any construction, describes the one architecture we know that passes every
test, and summarizes the mechanisms that fail. It contains no new theorems: every claim is either one of
the results above, a citation, or explicitly labelled as a heuristic.

\subsection{Necessary conditions}

A machine witnessing~(T) must have all of the following properties.
\begin{enumerate}[label=(N\arabic*),leftmargin=3em]
\item \emph{Coding.} It cannot only move data: machines that move indivisible items need time
$\Omega(L\log n)$ (Stoß~\cite{Stoss1973} for Turing machines, Aggarwal and Vitter~\cite{AggarwalVitter1988}
in external memory). Cells must hold functions of several entries.
\item \emph{Many passes, at every scale.} Its run has $\Omega(\log n)$ passes (Theorem~P$'$,
\cref{thm:Pprime}), even on a $2^{-O(nb)}$ fraction of the inputs (\cref{thm:Pprime-on}); the grouped form
(\cref{thm:Pgamma}) gives the same at coarser column scales.
\item \emph{Folding motion.} Along the popular trace, the macroscopic work-head trajectories satisfy $(\star)$
(\cref{cor:mstar-T}): some head makes a macroscopic U-turn, and the heads together waste at least one unit of
range (\cref{thm:excess}).
\item \emph{A maximal coding gain.} The block network of its run carries a causal code for complete exchange,
on a set of co-entropy $O(n)$, at the profile capacities (\cref{thm:reduction,thm:coent-linear}); by
\cref{thm:code-gap} this code beats routing by a factor $\Omega(\log n)$, the largest gain that cut bounds
allow. For an oblivious machine the code works on all inputs (\cref{thm:oblivious}); the same then holds on
the bit-level network~\cite[Thm.~6.2]{AHJKL2006}.
\item \emph{Not an XOR-frame recursion.} It cannot be built from the gadgets of \cref{sec:frames}, which cost at
least the trivial gadget.
\end{enumerate}

Condition~(N4) is the heart of the matter. Coding gains over routing in undirected networks do exist: the
multiple-unicast (network coding) conjecture is false~\cite{BravermanHe2026}. But the known growing gains are
of order $(\log n)^{\varepsilon}$, obtained by amplifying a constant gap by tensoring~\cite{BGS2017}, and they
concern sparse demands; tensoring a constant gap cannot produce the maximal gain $\Theta(\log n)$ for dense
demands (heuristically, and as observed by Zhang, Li and Li~\cite{ZhangLiLi2026}). A linear-time transposer
needs a \emph{native} maximal gain for the densest demand, on a network shaped by tapes.

\subsection{An architecture that survives}

We describe, as a design and not as a result, the one architecture that we could not rule out.

\paragraph{Geometric recursion.} The transposer runs $\Theta(\log n)$ levels. Level $\ell$ works on data of
volume $\theta^\ell L$ for a constant $\theta<1$, with $O(1)$ passes of an encoder $E_\ell$ and a decoder
$D_\ell$: it computes its output as $D_\ell(X_\ell,f_{\ell+1}(E_\ell(X_\ell)))$, where $f_{\ell+1}$ is the next
level's problem. The total time is $\sum_\ell O(\theta^\ell L)=O(L)$, and the pass count is $\Theta(\log n)$, as
Theorem~P requires.

\paragraph{Expanding choreography.} The heads of each level are driven along prepared \emph{rulers} that realize
two non-commuting monotone maps of the interval, so that the relation graph between the encoder's input and
output rows is a bounded-degree expander. Bourgain and Yehudayoff's monotone expanders~\cite{BourgainYehudayoff2013}
provide such maps; the rulers can be written in linear time, since their marks are computed with
$O(\log L)$-bit arithmetic. Near-identity generators are needed: integer or dyadic M\"obius pairs generate
discrete groups whose relation graphs are not expanders (a heuristic supported by numerical spectral gaps).
Repeating each choreography a constant number of times should make the macroscopic motion satisfy the
quantitative inequality $(\star)$; this is a cost constraint, not an obstruction (heuristic).

\paragraph{Two coding cores.} The geometry is the easy part. Each level needs a constant-fraction coding saving at
the encoder, so that $\abs{E_\ell(X_\ell)}\le\theta\abs{X_\ell}$, and, because the decoder writes the level's whole
output from $O(1)$-pass views of the next level's output, also at the decoder. In the simplest linear form the
encoder core is the following problem.

\begin{quote}
\emph{Locally encoded index coding with own-row side information.} Let $\Gamma$ be a bounded-degree graph on
$[n]$. A broadcast consists of linear functionals of an $n\times n$ matrix, each supported on the two rows of an
edge of $\Gamma$. Receiver $j$ sees the broadcast and knows row $j$, and must recover column $j$. Can the broadcast
have size $(1-\varepsilon)n^2$ for a fixed $\varepsilon>0$?
\end{quote}

Without the degree bound, $\varepsilon=\tfrac12$ is achievable: broadcast $A+A^{\T}$ above the diagonal. With the
degree bound, routing-type codes save only $O(n^2/\log n)$ on expanders, and hyperfinite (amenable) graphs save
$o(n^2)$ (heuristic arguments; not included here). A positive answer on expanders would be the first
constant-fraction coded saving in a tape-relevant model. A negative answer, together with its decoder analogue
(restricted-view index coding), would rule out every geometric recursion with $O(1)$-pass levels.

\subsection{Mechanisms that fail}

The following classes of constructions fail for reasons proved in this paper.
\begin{itemize}[leftmargin=2em]
\item \emph{Few passes.} Any schedule with $o(\log n)$ passes, with any coding, including constant-pass
choreographies along expanders (Theorems~P and~P$'$).
\item \emph{Regular or one-way motion.} Machines whose macroscopic head motion is stationary or monotone on every
tape (\cref{thm:excess,thm:mstar-adaptive}).
\item \emph{XOR-frame recursions,} including the template of~\cite{OpenAI2026mult}, which cannot reach linear time
(\cref{sec:frames}).
\item \emph{Codes whose linear structure covers the input set.} For a construction aimed at the plain conjecture's
counterexamples, linear codes, codes correct off a small set, and codes that are linear outside $O(nw)$ bits of
capacity reduce to plain codes (\cref{sec:plainrobust}).
\end{itemize}
Many further attempts failed for reasons we could only argue informally: flat schemes that code across a constant
number of whole streams (they save at most a constant factor), window-local recursions whose levels read only
nearby data, static coding certificates on Ramanujan-type graphs that resist conversion into causal codes, and
re-indexings that turn transposition into a cheap map only through a transposition-hard change of basis. We
mention them as warnings, not as results.
